\section{Introduction}
\label{sec:introduction}

% Paragraph 1 – broad motivation
Corporate disclosures contain both quantitative accounting information and qualitative narrative discussion that may convey incremental information about firm performance, risk, and prospects. A growing literature shows that markets respond not only to the numerical content of disclosures but also to their linguistic tone. However, in disclosure settings, economic inference depends importantly on how sentiment is measured. This issue is especially salient outside English-language contexts, where language structure, domain-specific usage, and translation frictions can all affect measurement reliability. This paper shows that, in Japanese MD\&A annual disclosures, sentiment measurement choice materially affects inference about filing-date market reactions, with a finance-specific Japanese lexicon exhibiting a stronger and more stable association with abnormal returns than a GPT-based score.

% Paragraph 2 – research gap / research question
These concerns are central in Japanese financial disclosures such as annual securities reports filed through EDINET. Most established financial sentiment tools were originally developed for English-language disclosures, and extending them to Japanese requires substantial linguistic and domain adaptation. At the same time, recent advances in large language models (LLMs), including GPT-based systems, suggest that contextual models may improve sentiment measurement by capturing meaning beyond fixed wordlists. This paper therefore asks whether sentiment extracted from Japanese MD\&A disclosures is reflected in stock price reactions around EDINET filing dates, and whether an LLM-based sentiment score provides incremental explanatory power beyond a finance-specific Japanese lexicon.

% Paragraph 3 – data
To answer this question, we use 1,087 MD\&A sections extracted from annual securities reports filed through EDINET by Nikkei 225 firms over 2018--2023. This setting is well suited to the research question because the JPFR 2017 taxonomy enables consistent identification of MD\&A text blocks, while the Nikkei 225 sample provides reliable market data and a large-cap institutional setting for studying disclosure reactions.

% Paragraph 4 – methods / empirical analysis
We construct two document-level sentiment measures. The first is a finance-specific Japanese adaptation of the Loughran--McDonald Master Dictionary, implemented with Japanese tokenization and normalization. The second is a GPT-based document sentiment score computed from the same MD\&A text. We then estimate cross-sectional event-study regressions of cumulative abnormal returns on standardized sentiment measures across multiple event windows, with year and industry fixed effects and firm-clustered standard errors.

% Paragraph 5 – main results
We find that lexicon-based sentiment is more consistently associated with abnormal returns than the GPT-based document score. GPT sentiment is positively related to returns in some standalone specifications, but its coefficient attenuates and becomes statistically indistinguishable from zero once Japanese LMMD tone is included, while LMMD remains economically meaningful in joint specifications. Thus, when abnormal returns are used as the benchmark for economically relevant disclosure tone, the Japanese LMMD measure exhibits a stronger and more persistent association across specifications and event windows in this setting.

% Paragraph 6 – contribution to literature
This paper contributes to three related literatures. First, it extends the literature on textual analysis of financial disclosures by constructing and validating a finance-specific Japanese sentiment measure for EDINET MD\&A filings. While prior work shows that narrative disclosure tone is informative in English-language settings \citep{tetlock2007giving,tetlock2008morethanwords,brown2011mda,feldman2010tone}, this paper studies a non-English regulatory environment in which language structure, financial terminology, and disclosure conventions create additional measurement challenges.

Second, it speaks to the literature on sentiment measurement in finance by comparing dictionary-based and contextual approaches in the same empirical setting. Prior studies emphasize the importance of domain-specific dictionaries for financial text \citep{loughran2011liability,henry2016measuring}, while recent work highlights the potential of LLM-based measures \citep{siano2025news,okada2025words}. This paper provides a direct comparison between a transparent finance-specific Japanese lexicon and a GPT-based sentiment score using the same disclosure sample, event dates, and stock-return outcomes.

Third, it informs the emerging literature on LLMs in empirical finance by showing that LLM-based sentiment need not dominate domain-specific lexical measures in structured regulatory filings. The results suggest that the value of LLM-based sentiment depends on the disclosure setting, language environment, and empirical objective.